% ===================================================================
%  Documentacao tecnica e didatica do pipeline de Tmax + Polymarket.
%  Os numeros do backtest vem de generated_numbers.tex, regenerado a
%  cada rodada do workflow de backtest — este documento nunca desatualiza.
%  Compile com: latexmk -pdf tmax.tex   (ou use Overleaf)
% ===================================================================
\documentclass[11pt,a4paper]{article}

\usepackage[utf8]{inputenc}
\usepackage[T1]{fontenc}
\usepackage[brazil]{babel}
\usepackage[margin=2.3cm]{geometry}
\usepackage{amsmath,amssymb}
\usepackage{booktabs}
\usepackage{xcolor}
\usepackage{fancyhdr}
\usepackage[colorlinks=true,linkcolor=blue!60!black,urlcolor=blue!60!black]{hyperref}

% Gerado por tmax/results.py a cada backtest — NÃO editar à mão.
% Atualizado em 2026-07-22T09:07:15+00:00

\newcommand{\edgeN}{57}
\newcommand{\edgeHit}{93\%}
\newcommand{\edgeComp}{2.97}
\newcommand{\edgeFlat}{+1.23}
\newcommand{\edgeDD}{2\%}
\newcommand{\edgeStops}{6}
\newcommand{\edgePrice}{0.82}
\newcommand{\combN}{381}
\newcommand{\combHit}{99\%}
\newcommand{\combComp}{3.77}
\newcommand{\combFlat}{+1.62}
\newcommand{\combDD}{2\%}
\newcommand{\combStops}{8}
\newcommand{\combPrice}{0.96}
\newcommand{\combMonthly}{+50\%}
\newcommand{\harvsixteenN}{324}
\newcommand{\harvsixteenHit}{100\%}
\newcommand{\harvsixteenComp}{1.40}
\newcommand{\harvsixteenFlat}{+0.39}
\newcommand{\harvsixteenDD}{1\%}
\newcommand{\harvsixteenStops}{2}
\newcommand{\harvsixteenPrice}{0.99}
\newcommand{\harvfourteenN}{603}
\newcommand{\harvfourteenHit}{100\%}
\newcommand{\harvfourteenComp}{1.63}
\newcommand{\harvfourteenFlat}{+0.66}
\newcommand{\harvfourteenDD}{5\%}
\newcommand{\harvfourteenStops}{11}
\newcommand{\harvfourteenPrice}{0.99}
\newcommand{\harvtwelveN}{823}
\newcommand{\harvtwelveHit}{100\%}
\newcommand{\harvtwelveComp}{1.87}
\newcommand{\harvtwelveFlat}{+0.94}
\newcommand{\harvtwelveDD}{6\%}
\newcommand{\harvtwelveStops}{17}
\newcommand{\harvtwelvePrice}{0.99}
\newcommand{\harvN}{324}
\newcommand{\harvHit}{100\%}
\newcommand{\harvComp}{1.40}
\newcommand{\harvFlat}{+0.39}
\newcommand{\harvDD}{1\%}
\newcommand{\harvStops}{2}
\newcommand{\harvPrice}{0.99}
\newcommand{\harvActiveHour}{16}
\newcommand{\confMin}{85\%}
\newcommand{\confMax}{93\%}
\newcommand{\btDaysCity}{1772}
\newcommand{\btSpan}{98}
\newcommand{\btStations}{18}
\newcommand{\btUpdated}{2026-07-22 09:07:15 UTC}


\definecolor{azul}{HTML}{1A5FB4}

% ---- ambiente de exercicios + resposta acumulada -----------------
\newcounter{exc}[section]
\renewcommand{\theexc}{\thesection.\arabic{exc}}
\newenvironment{exercicio}{\refstepcounter{exc}\medskip\noindent
  \textbf{\color{azul}Exercício \theexc.}\ }{\medskip}

\pagestyle{fancy}\fancyhf{}
\lhead{\small Trading Sistemático de Tmax}
\rhead{\small\thepage}
\renewcommand{\headrulewidth}{0.4pt}

\title{\textbf{Trading Sistemático de Temperatura Máxima}\\
  \large Documentação técnica e guia de estudo --- SBGR, SAEZ, UUWW e mais}
\author{Pipeline \texttt{tmax/} --- mercados de temperatura na Polymarket}
\date{Números do backtest atualizados em \btUpdated}

\begin{document}
\maketitle
\thispagestyle{fancy}

\begin{abstract}
\noindent Este documento descreve, com a matemática e a estatística por trás,
como o pipeline transforma previsões meteorológicas em decisões de trade em
mercados binários de temperatura máxima, e como validamos essas decisões com
backtests. É também um guia de estudo: cada seção termina com exercícios
(gabarito no apêndice). \textbf{Não é aconselhamento financeiro} --- mercados
de previsão envolvem risco de perda total do capital, e o \emph{edge} estimado
aqui é um modelo, não uma garantia.
\end{abstract}

\tableofcontents
\bigskip

% ===================================================================
\section{Visão geral}
% ===================================================================
O objetivo é comparar, para cada faixa de um mercado do tipo ``a máxima de hoje
vai ser $N^\circ$C?'', a nossa probabilidade estimada com a probabilidade
implícita pelo preço, e agir quando as duas divergem o suficiente e o modelo
está calibradamente confiante.

O sistema tem três camadas:
\begin{enumerate}
  \item \textbf{Previsão} (Seção~\ref{sec:previsao}): ensembles numéricos +
    correção de viés + observações em tempo real $\rightarrow$ uma
    \emph{distribuição de probabilidade} da máxima do dia.
  \item \textbf{Precificação e calibração} (Seções~\ref{sec:mercado}
    e~\ref{sec:calib}): a distribuição vira probabilidade por faixa de mercado,
    corrigida por uma curva de calibração empírica.
  \item \textbf{Estratégia e risco} (Seções~\ref{sec:estrategia}
    e~\ref{sec:risco}): quando comprar, quanto apostar, quando sair.
\end{enumerate}

Hoje o sistema é \emph{somente-leitura}: monitora \btStations{} cidades,
manda alertas ao Telegram e mede tudo com backtests; a ordem na Polymarket é
executada manualmente. Estações que resolvem em Fahrenheit foram catalogadas
mas ficam fora de operação.

\begin{exercicio}
Num binário, o preço do ``Sim'' (entre 0 e 1) é interpretado como a
probabilidade que o mercado atribui ao evento. Se o ``Sim'' custa \$0,40, qual
a probabilidade implícita do ``Não''?
\end{exercicio}

% ===================================================================
\section{Da previsão à distribuição}\label{sec:previsao}
% ===================================================================

\subsection{Por que ensemble}
Um modelo de tempo é um sistema caótico: erros minúsculos na condição inicial
crescem exponencialmente. Um \emph{ensemble} roda o mesmo modelo várias vezes
com perturbações na condição inicial (ECMWF ENS: 51 membros; GEFS: 31). Cada
membro é uma amostra do futuro possível; a fração de membros acima de um limiar
estima a probabilidade de excedê-lo. Usamos $\sim$82 membros.

\subsection{Correção de viés (\texttt{bias.py})}
Modelos globais têm erro sistemático num ponto específico. Estimamos, por
família de modelo $m$, ao longo dos últimos 60 dias:
\begin{align}
  \text{viés}_m &= \frac{1}{n}\sum_{d=1}^{n}\bigl[\hat T_m(d) - T_{\text{obs}}(d)\bigr],
    \qquad n \ge 15,\\
  \sigma^{\text{res}}_m &= \sqrt{\tfrac{1}{n-1}\sum_{d=1}^{n}
    \bigl[\hat T_m(d) - T_{\text{obs}}(d) - \text{viés}_m\bigr]^2 }.
\end{align}
A previsão corrigida é $\hat T_m - \text{viés}_m$; o desvio residual
$\sigma^{\text{res}}_m$ é o erro que \emph{sobra} e vira o $\sigma$ do membro na
mistura (Seção~\ref{sec:mistura}).

\subsection{Nowcast intradiário}
Durante o dia, o observado pesa mais que o modelo. Medimos o desvio médio entre
o METAR e a média do ensemble corrigido nas últimas 3 horas e deslocamos as
horas restantes por uma fração amortecida:
\begin{equation}
  \text{offset} = \operatorname*{mean}_{\text{últ. 3h}}
    \bigl[T_{\text{obs}}(h) - \bar T^{\text{corr}}_{\text{ens}}(h)\bigr],
  \qquad
  \text{shift} = 0{,}7 \cdot w(h)\cdot \text{offset},
\end{equation}
com $w(h)=\operatorname{clamp}\!\bigl(\tfrac{h-6}{6},\,0{,}25,\,1\bigr)$. O
fator $0{,}7$ evita perseguir ruído; $w(h)$ reduz o peso da madrugada, quando
nevoeiro e inversão térmica enganam.

\subsection{Mistura de gaussianas}\label{sec:mistura}
Para cada membro $i$, a máxima corrigida do dia é
$\hat T_i = \max(\text{obs\_max},\ \max_h[\text{horas futuras corrigidas}]+\text{shift})$.
A distribuição da máxima é a \emph{média} das CDFs (mistura com pesos iguais):
\begin{equation}
  F(x) = \Pr(T_{\max}\le x) = \frac{1}{M}\sum_{i=1}^{M}
    \Phi\!\left(\frac{x-\hat T_i}{\sigma_i}\right),
  \qquad
  \sigma_i = \max\!\bigl(\sigma^{\text{res}}_{\text{fam}(i)}\cdot f_{\text{rest}},\ 0{,}3\bigr),
\end{equation}
com $f_{\text{rest}}=\operatorname{clamp}(\tfrac{17-h}{12},0{,}2,1)$ ---
a incerteza encolhe à tarde. Duas cirurgias que os backtests motivaram:
\begin{itemize}
  \item \textbf{Truncagem:} $F(x)=0$ para $x<\text{obs\_max}$ --- a máxima nunca
    fica abaixo do já observado.
  \item \textbf{Trava:} quando o pico \emph{observado} já passou (3 horas
    seguidas abaixo da máxima) e nenhuma hora futura chega perto, $\sigma_i$
    colapsa para $0{,}05$. Exige a trava \emph{observacional} --- colapsar por
    previsão (o ensemble ``achando'' que a tarde não esquenta) gerava certezas
    falsas de manhã.
\end{itemize}
Os quantis (P10/P50/P90) saem por bissecção sobre $F$.

\begin{exercicio}
O viés do ECMWF em Guarulhos é $-0{,}7^\circ$C e o modelo prevê máxima de
$24{,}0^\circ$C. Qual a previsão corrigida? Se $\sigma^{\text{res}}=0{,}9$ e
$f_{\text{rest}}=1$, qual o $\sigma$ do membro?
\end{exercicio}
\begin{exercicio}
Mistura de dois membros, $\mathcal N(25,1)$ e $\mathcal N(27,1)$. Calcule
$\Pr(T_{\max}\le 26)$ usando $\Phi(1)=0{,}841$ e $\Phi(-1)=0{,}159$.
\end{exercicio}

% ===================================================================
\section{Do modelo ao mercado}\label{sec:mercado}
% ===================================================================
\subsection{Faixas, arredondamento e unidades}
O mercado resolve pela máxima \emph{arredondada ao inteiro} na unidade da
cidade. A faixa ``$N^\circ$C'' cobre $[N-0{,}5,\,N+0{,}5)$; ``entre $A$--$B\,^\circ$F''
cobre $[A-0{,}5,\,B+0{,}5)$ em Fahrenheit. Como a distribuição interna é em
$^\circ$C, faixas em $^\circ$F são convertidas nas \emph{bordas}:
$c(f)=(f-32)\cdot 5/9$. A probabilidade da faixa $[A,B]$ é
\begin{equation}
  \Pr(\text{Sim}) = F\!\bigl(c(B+0{,}5)\bigr) - F\!\bigl(c(A-0{,}5)\bigr).
\end{equation}

\subsection{Edge e valor esperado}
Num binário, o preço do Sim \emph{é} a probabilidade implícita $q$. Pagando $q$
num evento com probabilidade real $p$:
\begin{equation}
  \mathrm{EV}\text{ por \$1} = p\cdot\Bigl(\tfrac{1}{q}-1\Bigr) - (1-p)
    = \frac{p-q}{q}.
\end{equation}
O \emph{edge} em pontos percentuais é $p-q$. O EV relativo cresce quando $q$ é
pequeno --- por isso ``bilhetes de loteria'' baratos dominam o P\&L de backtests
ingênuos e merecem desconfiança (preço fino não executa).

\begin{exercicio}
Mercado NYC ``entre 78--79$^\circ$F''. Converta as bordas da faixa para
$^\circ$C.
\end{exercicio}
\begin{exercicio}
Você compra o ``Não'' a \$0,85 e sua probabilidade calibrada do ``Não'' é
$94\%$. Calcule o edge (p.p.) e o EV por \$1.
\end{exercicio}

% ===================================================================
\section{Calibração --- o coração estatístico}\label{sec:calib}
% ===================================================================
Um modelo é \emph{calibrado} se, entre os casos em que diz $90\%$, o evento
ocorre $\sim 90\%$ das vezes. Medimos com o \emph{Brier score} (erro quadrático
médio da probabilidade; menor é melhor):
\begin{equation}
  \text{Brier} = \frac{1}{n}\sum_{i=1}^{n}(p_i - y_i)^2, \qquad y_i\in\{0,1\}.
\end{equation}
Diagnóstico sobre $\sim$120 mil pares previsão--desfecho: o modelo cru declara
$\sim 99\%$ mas acerta entre \confMin{} e \confMax{} nas faixas-dia com
confiança $\ge 90\%$ (varia por cidade) --- superconfiança nas caudas.

\subsection{Regressão isotônica (PAVA)}
Ajustamos $g$ monótona não-decrescente minimizando $\sum(g(p_i)-y_i)^2$ pelo
algoritmo \emph{Pool Adjacent Violators}: percorre os pares ordenados por $p$ e
funde blocos vizinhos enquanto a média de um bloco anterior for $\ge$ à do
seguinte. Resultado: um mapa ``probabilidade dita $\to$ frequência real'', por
\emph{período do dia} (0--5h / 6--11h / 12--23h), porque a superconfiança da
madrugada (pouca observação) é muito maior que a da tarde.

\subsection{O \emph{blend} com o preço (diagnóstico)}
Uma regressão logística $y\sim a\,\operatorname{logit}(p_{\text{modelo}})
+ b\,\operatorname{logit}(\text{preço})+c$, por período, mede quem carrega
informação: de madrugada $a\approx 0$ (o modelo não adiciona nada ao preço);
à tarde $a>0$ (adiciona). Decisão operacional: os sinais usam só o modelo
calibrado (o mesmo número da tabela); o blend segue como diagnóstico.

\begin{exercicio}
O modelo declara $99\%$ em 842 faixas-dia e acerta $86\%$. Qual o Brier
aproximado dessas previsões? (aproxime $p=0{,}99$ e frequência $0{,}86$.)
\end{exercicio}
\begin{exercicio}
PAVA: pares ordenados por $p$ com desfechos $y=[0,1,0,1,1]$. Quais os blocos e
os valores finais da isotônica?
\end{exercicio}

% ===================================================================
\section{Estratégias operacionais}\label{sec:estrategia}
% ===================================================================
\subsection{Estratégia \emph{Edge}}
Compra o ``Não'' de uma faixa de hoje quando:
\begin{itemize}
  \item $|p_{\text{cal}} - q| \ge 5$ p.p. (divergência mínima);
  \item o lado ``Não'' tem $>90\%$ de chance segundo o modelo calibrado;
  \item preço do ``Não'' $\ge \$0{,}30$ (não brigar com mercado quase-certo);
  \item hora local entre 6h e 23h; cruzamento novo (madrugada é descartada).
\end{itemize}
Só o lado ``Não'' é operado: o ``Sim'' teve $\sim 41\%$ de acerto (loteria).
Testamos e \emph{rejeitamos} um teto de edge de 20 p.p.: no lado ``Não'' já
filtrado, gaps grandes são o modelo vendo a máxima da tarde antes do mercado
convergir, não erro.

\subsection{Estratégia \emph{Colheita de favoritos}}
Compra o ``Não'' quase-certo (preço entre \$0,97 e \$0,995), após as
\harvActiveHour{}h locais, com o modelo concordando. Depois desse horário a
máxima está praticamente travada e o favorito a 97 centavos é quase-matemática.
Colhe os últimos centavos de incerteza residual.

\subsection{Stop loss}
Num binário o preço \emph{é} a probabilidade: cair $15\%$ significa que o
mercado agora acha o evento $15\%$ menos provável. O stop converte a perda
potencial total ($-100\%$) numa perda pequena e certa ($-15\%$), ao custo de
realizar ruído. Vale quando o gatilho é \emph{informativo} (o dado virou --- a
máxima estourou a faixa) e custa caro quando é só volatilidade.

\begin{exercicio}
Você entra a \$0,70; o alerta dispara a $-10\%$ e a saída do backtest a
$-15\%$. Em que preços cada um dispara? Se o evento venceria na resolução,
quanto o stop custou por \$1 de stake?
\end{exercicio}

\subsection{Como os alertas chegam ao vivo (produção)}\label{sec:disparo-vivo}
O backtest conta \emph{entradas} (uma por faixa/dia, no primeiro cruzamento); o
que chega ao Telegram segue regras de \emph{entrega} diferentes:
\begin{itemize}
  \item \textbf{Modo silencioso:} o chat recebe só alertas (compra, colheita,
  condição, stop) e o resumo de posições no máximo $1\times$ por hora.
  \item \textbf{Alerta de compra novo:} chega com o \emph{bloco completo} da
  cidade (tabela, gráfico, hora a hora) --- o contexto da decisão de entrada.
  Você também pode pedir o relatório de qualquer cidade a qualquer momento por
  \texttt{/relatorio}.
  \item \textbf{Repetição:} \emph{edge} e colheita repetem a cada rodada
  ($\sim$10\,min) enquanto a oportunidade durar; as repetições vêm sozinhas em
  texto curto (sem repetir os gráficos). Uma entrada equivale a muitas
  mensagens.
  \item \textbf{Alertas de condição} (só cidades com posição aberta): platô de
  2h de lado e fuga do envelope do \emph{ensemble} --- uma vez por episódio.
  \item \textbf{Stop:} dispara a $-10\%$ do preço de entrada e repete toda
  rodada; a saída de referência do backtest é $-15\%$.
\end{itemize}

\paragraph{Ao vivo dispara mais que o backtest.} O backtest só considera faixas
com negócio recente ($\le$2h); ao vivo o alerta usa o preço nominal exibido, e
por isso também dispara em faixas ilíquidas que o backtest descarta. Medido no
arquivo atual, são $\sim$18\% mais sinais ao vivo do que a simulação reporta ---
não é erro, são medições diferentes. Além disso, num único dia os sinais se
\emph{acumulam} no chat (cada cidade no seu fim de tarde local), enquanto o
backtest os distribui por data: por isso um dia com 3--4 sinais parece muito,
mas ainda é um evento raro (o histórico teve no máximo 4 num único dia).

% ===================================================================
\section{Backtest --- metodologia e resultados atuais}\label{sec:backtest}
% ===================================================================
\subsection{Metodologia}
Para cada dia arquivado reconstruímos, hora a hora, o que o sistema atual teria
dito: viés ``como era'' (janela de 60 dias terminando 2 dias antes), METARs até
a hora $h$, ensemble arquivado, nowcast, mistura $\rightarrow$ probabilidade por
faixa; preço $=$ último negócio (máx.\ 2h de idade); resolução $=$ desfecho
oficial. Aposta de $10\%$ do capital no primeiro cruzamento da regra, liquidada
na resolução (ou no stop de $-15\%$). \textbf{Flat} usa sempre $10\%$ do capital
\emph{inicial} (para comparar regras); \textbf{composto} reinveste $10\%$ do
capital \emph{corrente} (o ``quanto eu teria'').

\paragraph{Vieses conhecidos.} \emph{Lookahead} do ensemble arquivado
(otimista); \emph{liquidez} (o último preço de mercado fino não executa em
tamanho); \emph{in-sample} (a calibração é ajustada no mesmo período --- mitigada
pelo reajuste a cada 3 dias, com dias novos entrando \emph{depois} do ajuste).

\subsection{Resultados (arquivo de \btDaysCity{} dias-cidade, $\sim$\btSpan{} dias corridos)}
\begin{center}
\begin{tabular}{lrrrrr}
\toprule
\textbf{Estratégia} & \textbf{Entradas} & \textbf{Acerto} & \textbf{Composto}
  & \textbf{DD máx} & \textbf{Stops}\\
\midrule
Edge (Não, simulação histórica) & \edgeN & \edgeHit & \edgeComp\,x & \edgeDD & \edgeStops\\
Colheita 16h & \harvsixteenN & \harvsixteenHit & \harvsixteenComp\,x & \harvsixteenDD & \harvsixteenStops\\
Colheita 14h & \harvfourteenN & \harvfourteenHit & \harvfourteenComp\,x & \harvfourteenDD & \harvfourteenStops\\
Colheita 12h & \harvtwelveN & \harvtwelveHit & \harvtwelveComp\,x & \harvtwelveDD & \harvtwelveStops\\
\midrule
\textbf{Combinado (Edge + Colheita \harvActiveHour h)} & \textbf{\combN} & \textbf{\combHit}
  & \textbf{\combComp\,x} & \textbf{\combDD} & \textbf{\combStops}\\
\bottomrule
\end{tabular}
\end{center}

\noindent Retorno composto do portfólio combinado, anualizando a janela de
$\sim$\btSpan{} dias para base mensal: \textbf{\combMonthly{} ao mês}
(in-sample; a expectativa executada realista é mais modesta, tipicamente
$10$--$20\%$/mês, por causa de liquidez e slippage). Preço médio pago pela
Edge: \$\edgePrice.

\emph{Estes números são regenerados automaticamente a cada backtest
(\btUpdated) --- veja \texttt{backtest\_results/} e
\texttt{generated\_numbers.tex}.}

\begin{exercicio}
A estratégia Edge fez \edgeN{} entradas em $\sim$\btSpan{} dias. Isso é o
número de alertas que você recebeu? Explique a diferença entre ``entradas
simuladas'' e ``mensagens recebidas''.
\end{exercicio}

% ===================================================================
\section{Gestão de risco}\label{sec:risco}
% ===================================================================
\emph{Kelly} para um binário comprado a $q$ com probabilidade real $p$:
\begin{equation}
  f^\star = \frac{p-q}{1-q}\quad(\text{fração da banca; aposte 0 se }p\le q).
\end{equation}
Kelly maximiza o crescimento logarítmico \emph{mas} assume $p$ exato --- e a
Seção~\ref{sec:calib} mostrou que nosso $p$ tem erro. Regra da casa: $\tfrac14$
de Kelly, teto de $5\%$ por mercado, $10$--$15\%$ somando o dia/cidade (faixas
do mesmo dia são \emph{uma} aposta), e \emph{circuit breaker} de drawdown
mensal.

\begin{exercicio}
Kelly: $p=0{,}90$, preço $q=0{,}80$. Calcule $f^\star$ e $\tfrac14$ de Kelly.
Com banca de \$2.000, quanto apostar (respeitando o teto de $5\%$)?
\end{exercicio}

% ===================================================================
\appendix
\section{Gabarito}
\begin{enumerate}
  \item[\textbf{1.1}] $1-0{,}40=0{,}60=60\%$.
  \item[\textbf{2.1}] $24{,}0-(-0{,}7)=24{,}7^\circ$C; $\sigma=\max(0{,}9\cdot1;\,0{,}3)=0{,}9$.
  \item[\textbf{2.2}] $F(26)=\tfrac12[\Phi(1)+\Phi(-1)]=\tfrac12(0{,}841+0{,}159)=0{,}50$.
  \item[\textbf{3.1}] $c(77{,}5)=(77{,}5-32)\cdot5/9=25{,}28^\circ$C;\ $c(79{,}5)=26{,}39^\circ$C. Faixa $[25{,}28;\,26{,}39)$.
  \item[\textbf{3.2}] Edge $=94-85=9$ p.p.;\ EV $=(0{,}94-0{,}85)/0{,}85=+10{,}6\%$ por \$1.
  \item[\textbf{4.1}] $\text{Brier}\approx 0{,}86\cdot(0{,}99-1)^2 + 0{,}14\cdot(0{,}99-0)^2 \approx 0{,}137$ (contra $\sim0{,}01$ se calibrado).
  \item[\textbf{4.2}] $[1,0]$ viola $\to$ funde em $0{,}5$. Valores finais: $0;\ 0{,}5;\ 0{,}5;\ 1;\ 1$.
  \item[\textbf{5.1}] Alerta a $0{,}70\cdot0{,}90=\$0{,}63$; saída a $0{,}70\cdot0{,}85=\$0{,}595$. Num vencedor você recebeu $0{,}85$ em vez de $1/0{,}70=1{,}43$ por \$1 de stake $\Rightarrow$ deixou $\sim0{,}58$ na mesa.
  \item[\textbf{6.1}] NÃO é contagem de alertas. ``Entradas simuladas'' $=$ quantas vezes a regra teria comprado reencenando o histórico; ``mensagens recebidas'' inclui as repetições a cada 10 min enquanto a oportunidade dura (uma oportunidade $=$ muitas mensagens).
  \item[\textbf{7.1}] $f^\star=(0{,}90-0{,}80)/(1-0{,}80)=0{,}50$;\ $\tfrac14$ Kelly $=12{,}5\%\Rightarrow\$250$, mas o teto de $5\%$ limita a \$100.
\end{enumerate}

\vfill
\noindent\rule{\linewidth}{0.4pt}\\
{\small Repositório \texttt{Temperatura-M-xima-de-CP}. Os números da
Seção~\ref{sec:backtest} vêm da última rodada do backtest e mudam conforme o
arquivo cresce. Documento gerado para estudo; revise as fontes antes de operar.}

\end{document}
